% p2-05-symbolic-rg

\section{P2 §5. Symbolic Renormalization Group: A Toy Model}

5. Symbolic Renormalization Group: A Toy Model

   Figure (fig-p2-ch5-symbolic-rg). Symbolic RG triptych --- Symbolic coupling space / beta-function
                                with zero / fixed-point eigenvalue ratio.

  [CONJECTURAL --- this section contains no derived results; it develops a formal
  language for later work packages]

  5.1 Motivation for a Symbolic RG
  Standard RG analysis operates on real-valued couplings. The results of Paper 1 suggest
  an alternative representation: couplings encoded as symbolic expressions over an
  algebraic basis. The symbolic RG (sRG) model asks whether RG flows act simply on
  such symbolic representations. The Lucas ring Z[$\varphi$] provides a natural candidate for the
  symbolic coupling space because it is closed under the $\varphi$2 = $\varphi$+1 recursion and carries
  the Coq-proven identity (1).

  5.2 Definition of the Symbolic Coupling Space [L3 --- Model-derived]
  Let Sigma be the set of algebraic numbers expressible as

  Sigma = { sumk ak $\varphi$nk : ak $\in$ Q, nk $\in$ Z, finitely many terms } = Q($\varphi$) = Q(5). 12

  A $\varphi$-structured coupling at scale mu is a coupling g(mu) whose minimal description
  length in a $\varphi$-basis is shorter than in a generic algebraic basis.

  The symbolic complexity C[g] of a coupling is defined operationally as the length of the
  shortest expression tree in a basis B = {$\varphi$, pi, e, +, -, $\times$, /} that evaluates to g within a
  specified numerical tolerance epsilon. Paper 1's finding is that for a selected set of
  coupling combinations, CB[g] < CB'[g] where B' replaces $\varphi$ with a generic algebraic
  constant. The robustness of this finding is Work Package 5.

  5.3 Symbolic Beta Functions [CONJECTURAL --- L5]
  A symbolic beta function is a map

  beta: Sigma $\rightarrow$ Sigma, g ↦ mu (dg)/(dmu) 13

  such that beta(g) $\in$ Sigma whenever g $\in$ Sigma. A $\varphi$-closed RG trajectory is a flow
  whose couplings remain in Sigma for all scales mu. The question of whether any
  realistic QFT beta function is $\varphi$-closed is open.

  The one-loop beta function for a single scalar coupling lambda in phi4 theory in 4D is
  (Wilson and Kogut, Phys. Rep. 12 (1974) 75):

  beta(lambda) = (3lambda2)/(16pi2) + O(lambda3). 14

  [Structural obstacle --- L3] At one loop, beta(lambda) $\in$ Sigma only if lambda $\in$ Sigma
  and 16pi2 $\in$ Sigma. Since pi2 notin Q(5) --- these are algebraically independent over Q
  by Nesterenko's theorem --- the naive $\varphi$-closure fails immediately at one loop. This is
  the $\pi^{2}$-incommensurability obstacle and is a concrete structural result, not a conjecture.
  It implies that $\varphi$-closure cannot hold for individual coupling beta functions at standard
  perturbative order; it must emerge in ratios or at fixed points, if it emerges at all.

  Possible resolutions [L5 --- all conjectural]:

  - $\varphi$-structure is approximate and emerges only in specific running-coupling ratios, not
  individual couplings.

  - The relevant quantities are operator eigenvalue ratios at fixed points, not the
  couplings themselves; these are pure numbers that may lie in Q(5).
  - $\varphi$-closure holds only at discrete RG scales separated by a ratio of $\varphi$, i.e., a DSI
  condition.
  - SUSY non-renormalization theorems protect certain operator dimensions exactly; a
  SUSY theory with a $\varphi$-eigenvalue ratio at tree level would maintain it to all orders.

  5.4 Fixed Points and Symbolic Eigenvalue Ratios [CONJECTURAL ---
  L5]
  At a fixed point g\textasciicircum{}*, the stability matrix

  Mij = ($\partial$ betai)/($\partial$ gj)|g\textasciicircum{}* 15

  has eigenvalues Deltai (the anomalous dimensions of critical operators). The ratio
  Deltaj / Deltai is a dimensionless number. The $\varphi$-Eigenvalue Conjecture reads:

          $\varphi$-Eigenvalue Conjecture [L5]. There exist physically motivated field theories in
          which Deltaj / Deltai = $\varphi$n for integer n at an RG fixed point.

  This is stated as a conjecture motivated by the Toda/Coxeter examples but not derived
  from them. A demonstration would require identifying an explicit fixed point with a
  $\varphi$-eigenvalue ratio and verifying it survives quantum corrections. This is Work Package
  1. The formal falsification criterion is given in the table in Section 4.4.

  The Trinity anchor identity is relevant here as a consistency check: if two scaling
  dimensions Delta1, Delta2 are in the ratio Delta2/Delta1 = $\varphi$, then Delta22/Delta12 =
  $\varphi$2 = $\varphi$ + 1. The identity $\varphi$2 + $\varphi$-2 = 3 constrains Delta1-2 + Delta2-2 = Delta1-2(1 +
  $\varphi$-2) and Delta12 + Delta22 = Delta12(1 + $\varphi$2) = Delta12(2 + $\varphi$). Neither expression
  equals 3 in general; the anchor constrains a specific bilinear form in $\varphi$ and $\varphi$-1, not
  arbitrary scaling dimension sums.

\subsection{References}

- Coldea, R. et al. "Quantum Criticality in an Ising Chain: Experimental Evidence for Emergent E8 Symmetry." \textit{Science} 327 (2010) 177. DOI \href{https://doi.org/10.1126/science.1180085}{10.1126/science.1180085}.
- Wu, J. et al. "E8 Symmetry in CoN$b_{2}$$O_{6}$ at Higher Resolution." \textit{Phys. Rev. B} 108, L060403 (2023). DOI \href{https://doi.org/10.1103/PhysRevB.108.L060403}{10.1103/PhysRevB.108.L060403}.
- Fring, A.; Korff, C. "Affine Toda field theories related to Coxeter groups of non-crystallographic type." \textit{Nucl. Phys. B} 729 (2005) 361. \href{https://arxiv.org/abs/hep-th/0506226}{arXiv:hep-th/0506226}. DOI \href{https://doi.org/10.1016/j.nuclphysb.2005.08.044}{10.1016/j.nuclphysb.2005.08.044}.
- Pellis, S. "The Fine-Structure Constant and the Golden Ratio." \href{https://vixra.org/pdf/2110.0117v4.pdf}{viXra 2110.0117 v4}.
